% answers_short.tex —— 2008 年·普通高等学校招生考试·四川卷·理科 · 简短答案速查
% 由 tools/make_short_answers.py 自 answers.tex 的 \gksolhead/\gkanswer 提取，
% 并与 2008/2008_数学真题_四川理_答案要点.md 逐题比对后生成，勿手改（改源 MD 后重跑）。

\begin{gknotice}
\textbf{使用说明}：本节只给最终答案便于快速核对，完整推导见解析版。答案为编写者独立推导，\textbf{非官方标准答案}。
\end{gknotice}

\gkpart{选择题}
\begin{center}\begin{tabular}{ccccccccccccc}
\toprule
题号 & 1 & 2 & 3 & 4 & 5 & 6 & 7 & 8 & 9 & 10 & 11 & 12 \\
\midrule
答案 & B & A & D & A & B & C & D & D & B & D & C & B \\
\bottomrule
\end{tabular}\end{center}

\gkpart{填空题}
\begin{itemize}[leftmargin=2.6em]
  \item[\textbf{13.}] $-6$。
  \item[\textbf{14.}] $\sqrt2$。
  \item[\textbf{15.}] $2$。
  \item[\textbf{16.}] $4$。
\end{itemize}

\gkpart{解答题}
\begin{itemize}[leftmargin=2.6em]
  \item[\textbf{17.}] 最大值 $10$，最小值 $6$。
  \item[\textbf{18.}] （1）$0.5$；（2）$0.8$；（3）分布列见解析，$E\xi=2.4$。
  \item[\textbf{19.}] （1）证明见解析；（2）$\arccos\dfrac{\sqrt{10}}5$。
  \item[\textbf{20.}] （1）证明见解析；（2）$b=2$ 时 $a_n=(n+1)2^{n-1}$。
  \item[\textbf{21.}] （1）$a=2$，$b=\sqrt2$；（2）证明见解析。
  \item[\textbf{22.}] （1）$a=16$；（2）增区间 $(-1,1)$、$(3,+\infty)$，减区间 $(1,3)$；（3）$\left(32\ln2-21,\ 16\ln2-9\right)$。
\end{itemize}
